\phantomsection
\chapter*{Kitabı Kullanma Rehberi}
\markboth{Kitabı Kullanma Rehberi}{Kitabı Kullanma Rehberi}
\addcontentsline{toc}{chapter}{Kitabı Kullanma Rehberi}

Kitabı araştırma sorusundan gerekçeli rapora uzanan bir çalışma rotası
olarak kullanın. Önce hangi soruyu yanıtlamak istediğinizi ve verinin nasıl
üretildiğini belirleyin. Ardından uygun özeti veya analizi seçin;
belirsizliği ve yorumun sınırlarını açıklayın. Bir yazılımın sonuç üretmesi,
bu adımların tamamlandığı anlamına gelmez.

\textbf{Ortak bölüm düzeni.} Birleşik kitabın 20 bölümündeki
\emph{Literatür verisiyle yazılım uygulamaları}, aynı verinin IBM SPSS,
R ve Python ile analizini, sonuç karşılaştırmasını ve örnek raporunu içerir.
Numaralı kutular SPSS menü adımlarını gösterir. Ayrı
\emph{R ile çözümlü uygulama} bölümleri ve konu anlatımındaki Python
örnekleri korunmuştur; bunların öğretim verileri ortak gerçek veriyle
karıştırılmamalıdır. Öğrenciden üç yazılımı aynı anda öğrenmesi beklenmez.
Bootstrap ve rastgeleleştirme, eksik veri, ölçek güvenirliği ve çoklu
regresyon bölümleri gerektiğinde destek ya da ileri okuma rotası sunar.

\begin{prismbox}[PrismLight]{Önerilen haftalık çalışma döngüsü}
\begin{enumerate}
  \item Bölümün amaç kutusunu ve ana kavramları okuyun.
  \item Araştırma sorusunu, gözlem birimini ve hesaplanacak büyüklüğü yazın; sonra formüldeki simgeleri açıklayın.
  \item Çözümlü örneği kapatarak yeniden çözün ve sonucu bir cümleyle yorumlayın.
  \item Yazılım çıktısını elle yapılan hesap ve grafiklerle karşılaştırın.
  \item Yöntem seçiminin gerekçesini, sonucun belirsizliğini ve ileri süremeyeceğiniz bir iddiayı yazın.
\end{enumerate}
\end{prismbox}

Kitaptaki \emph{Sezgi} kutuları kavramsal bağlantıları, \emph{Sık Hata}
kutuları yaygın yanlış yorumları, \emph{Yazılım Notu} ve
\emph{Yazılım Çıktısı Nasıl Okunur?} kutuları ise teknoloji destekli analizde
denetlenmesi gereken noktaları gösterir. Bu kutular ana metnin süsü değil,
istatistiksel karar sürecinin parçasıdır. Bu kullanım düzeni, kavramsal
anlama ile gerçek veri çözümlemesini birlikte geliştiren önerilerle uyumludur
\citep{gaise2016}.

\textbf{Ortak pedagojik omurga.} Her ana bölüm; öğrenme hedefleri, başlamadan
önce sorusu veya ön koşul bilgisi, temel kavram açıklaması, en az bir çözümlü
uygulama, yöntem/çıktı denetimi, bölüm özeti ve uygulama soruları içerir.
Bölümün niteliğine göre gerçek veri laboratuvarı ile Python, SPSS veya R
uygulamalarından uygun olanlar eklenir; üç yazılımın her bölümde zorunlu ve
eşit uzunlukta kullanılması beklenmez.

Hesap makinesi veya yazılım aritmetik yükünü azaltabilir; fakat hedef
parametreyi, uygun yöntemi ve yorumun sınırlarını kendi başına belirleyemez.
Bu nedenle her analizde sayısal sonuçtan önce araştırma tasarımı, sonra
varsayımlar ve en son bağlamsal yorum ele alınmalıdır.

\section*{Karar ve raporlama kaydı}

Aşağıdaki kaydı uygulamalara eşlik eden kendi çalışma notunuz olarak
kullanın. Bu, her bölümde doldurulmuş olarak bulunan bir tablo değil,
farklı örneklere uygulayabileceğiniz ortak bir kontrol listesidir.
Yalnız betimleme yapılan bir görevde test veya güven aralığı gerekmiyorsa,
bunları zorunlu olarak eklemek yerine neden gerekmediğini belirtin.

\begin{enumerate}
  \item \textbf{Soru ve veri:} Araştırma sorusu nedir? Gözlem birimi,
  veri kaynağı ve araştırma tasarımı nedir; hangi kayıtlar kullanılmıştır?
  \item \textbf{Hedef:} Ortalama, oran, fark veya ilişki gibi hangi
  büyüklük incelenmektedir? Örneklem özeti ile evrene ilişkin hedefi ayırın.
  \item \textbf{Analiz tanımı:} Yöntemi, değişken veya ölçüm sırasını,
  eksik değer yaklaşımını ve varsa test yönü ile güven düzeyini yazın.
  Yöntemin koşullarını ve değerlendiremediğiniz varsayımları belirtin.
  \item \textbf{Sonuç ve belirsizlik:} Tahmini birimiyle verin; soruya
  uygun olduğunda güven aralığını, ilişki veya fark büyüklüğünü ve test
  sonucunu ekleyin. Yalnız $p$ değeriyle yetinmeyin.
  \item \textbf{Yorumun sınırı:} Soruyu yanıtlayan kısa bir rapor yazın.
  Verinin desteklemediği bir nedensellik, genelleme veya tahmin iddiasını
  ayrıca belirtin.
\end{enumerate}

\textbf{Aynı veri, farklı soru.} B11'deki gruplararası ortalama farkı ile
B12'deki 10 puan eşiğini karşılama oranlarını karşılaştırırken önce hangi
büyüklüğün değiştiğini açıklayın. B14'te yalnız pozitif notları seçmenin
hangi kayıtları ve hedefi değiştirdiğini düşünün. Farklı sonuçları yalnız
``anlamlı/anlamsız'' diye ayırmak yerine, her birinin yanıtladığı soruyu
yazın. Aynı kayıtları yeniden analiz etmek bağımsız bir tekrar çalışma
değildir.

\textbf{Çıktıları eşleştirme.} Yazılımlar arasında karşılaştırmaya geçmeden
önce analiz tanımının aynı olduğunu kontrol edin. Farkın yönü, testin iki
yönlü olması veya süreklilik düzeltmesi gibi tercihler sonuçları
etkileyebilir. Kullandığınız kodu ve çıktıyı saklayın; yeniden çalıştırma
ile yalnız paylaşılan bir çıktıyı incelemeyi birbirinden ayırın.
Sayısal uyum, veri üretiminin veya model varsayımlarının denetlendiği
anlamına gelmez.

Hızlı tekrar sırasında formüller için Ek~\ref{app:formulas}, kritik değerler
için Ek~\ref{app:distribution-tables}, yazılım eşleştirmeleri için
Ek~\ref{app:software-map} kullanılabilir. Bölüm sonu sorularında önce yalnız
ipucunu, ardından kısa yanıtı görmek isteyen okuyucu Ek~\ref{app:answer-key}'deki
kademeli anahtara başvurabilir.

\section*{Gerçek veri ve kod paketi}

Kitabın ortak veri ve kod dosyaları \texttt{companion/} klasöründe bulunur.
Bu klasörde ham ve temiz veri ayrı tutulmuş; değişken sözlükleri, veri
kaynakları, bölüm eşleştirmeleri ve yeniden üretilebilir Python betikleri
birlikte verilmiştir. \texttt{companion/code/} içindeki veri hazırlama,
doğrulama, örnek analiz ve grafik üretim akışı proje kökünden
\begin{center}
\ifimoders
\texttt{python companion/code/run\_imo301.py}
\else
\texttt{python companion/code/run\_all.py}
\fi
\end{center}
komutuyla çalıştırılabilir. Bu komut bütün basılı yazılım bloklarını veya
SPSS ve R'yi çalıştırmaz. Literatür uygulamalarının kodları için ilgili
bölümü ve \texttt{companion/spss/} paketini izleyin; hazırlık adımları
sayfa~\pageref{sec:spss-guide}'de açıklanmıştır.
\ifimoders
Ortak paket kapsamlı kitabın dosyalarını da içerir; bu komut yalnız ders
sürümünün temel analizlerini çalıştırır. Çıktı dosyalarındaki numaralar
kaynak dosya adlarıyla eşleşir ve bu sürümdeki bölüm sırasını izler.
\fi

Gerçek veri setleri kullanılırken gözlem birimi ve veri kaynağı her analizde
yeniden belirtilmelidir. Okul bölgesi düzeyindeki bir ilişki, öğrenci
düzeyindeki ilişki gibi yorumlanmamalıdır. Veri setlerinin kaynak ve hak
bilgileri \texttt{companion/data/dictionaries/} altında kayıtlıdır.
Verinin kaynağını, kullanım koşullarını ve yapılan dönüşümleri belgelemek,
istatistiksel uygulamanın etik sorumlulukları arasındadır
\citep{asaethics2022}.

\begin{mistakebox}
Bir formülü doğru uygulamak, yanlış seçilmiş örneklemi veya uygun olmayan
araştırma tasarımını düzeltmez. İstatistiksel kalite veri üretim aşamasında
başlar.
\end{mistakebox}